\documentclass[11pt,a4paper]{ctexart}
\input{/Users/shaoyikai/Desktop/考研/22408-notes-skill/assets/note-preamble.tex}
\renewcommand{\NoteFooter}{22408 一页纸笔记 · 数学二}
\linespread{1.14}
\begin{document}
\notetitle{旋转体体积与表面积}{数学二}{高等数学·定积分应用}
\sechead{一、绕 $x$ 轴旋转的核心公式}

设曲线光滑（或分段光滑）；一般情形下侧面积公式中的 $y$ 需取绝对值 $\lvert y\rvert$：

\begin{center}\small
\begin{tabularx}{\linewidth}{@{}>{\raggedright\arraybackslash}X>{\raggedright\arraybackslash}X>{\raggedright\arraybackslash}X>{\raggedright\arraybackslash}X@{}}
\toprule
曲线表示 & 旋转体体积 $V$ & 旋转体侧面积 $S$ & 适用条件 \\
\midrule
直角坐标 $y=f(x)$ & $\pi\int_a^b y^{2}\,\mathrm{d}x$ & $2\pi\int_a^b \lvert y\rvert\sqrt{1+y'^{2}}\,\mathrm{d}x$ & $a\le x\le b$，连续/分段光滑 \\
参数方程 $x(t),y(t)$ & $\pi\int_\alpha^\beta y^{2}(t)\lvert x'(t)\rvert\,\mathrm{d}t$ & $2\pi\int_\alpha^\beta \lvert y(t)\rvert\sqrt{x'^{2}+y'^{2}}\,\mathrm{d}t$ & $\alpha\le t\le\beta$，$x(t)$ 单调 \\
双边界（垫圈法） & $\pi\int_a^b (y_{\text{上}}^{2}-y_{\text{下}}^{2})\,\mathrm{d}x$ & $S_{\text{上}}+S_{\text{下}}$（外、内侧面之和） & $0\le y_{\text{下}}(x)\le y_{\text{上}}(x)$ \\
\bottomrule
\end{tabularx}
\end{center}

\notebox{参数方程体积定限技巧：若按 $x$ 从 $a$ 增到 $b$ 积分，让 $t$ 从 $x=a$ 对应的 $t_1$ 积到 $x=b$ 对应的 $t_2$，直接写 $\displaystyle\pi\int_{t_1}^{t_2} y^{2}(t)x'(t)\,\mathrm{d}t$ 即可（符号自动抵消）。}

\sechead{二、两曲线围成区域的“加减法则”}

设区域 $D$ 位于 $x$ 轴上方（$0\le y_{\text{下}}(x)\le y_{\text{上}}(x)$，$a\le x\le b$），由上下两曲线围成：

\begin{itemize}

\item \textbf{判上下}：先求交点确定 $[a,b]$，再比较同 $x$ 处（利用不等式或单调性）谁的 $y$ 更大。

\item \textbf{体积相减（挖空）}：外层旋转体减去内层空腔，即 $V = V_{\text{上}} - V_{\text{下}}$。

\item \textbf{表面积相加（内外表面）}：空腔内壁也是外露表面，故 $S = S_{\text{上}} + S_{\text{下}}$；若端点另有竖直截线段 $x=a$ 或 $x=b$，再加圆环端面面积 $\pi(y_{\text{上}}^2 - y_{\text{下}}^2)$。

\end{itemize}

\sechead{三、典型例题：圆弧与星形线}

\textbf{例} 设 $D$ 由 $y=\sqrt{1-x^{2}}\ (0\le x\le 1)$ 与星形线 $x=\cos^{3}t,\ y=\sin^{3}t\ (0\le t\le\frac{\pi}{2})$ 围成，求 $D$ 绕 $x$ 轴旋转一周所得旋转体的体积 $V$ 和表面积 $S$。

\textbf{1. 定边界与上下关系}：两曲线均连接 $(0,1)$ 与 $(1,0)$，无竖直端面。对 $0<t<\frac{\pi}{2}$，有 $x,y\in(0,1)$，故 $x^{2}+y^{2} < x^{2/3}+y^{2/3}=1$，即\textbf{圆弧在上、星形线在下}。

\textbf{2. 求体积 $V = V_{\text{圆}} - V_{\text{星}}$}：

\[
V_{\text{圆}}=\pi\int_{0}^{1}(1-x^{2})\,\mathrm{d}x=\frac{2}{3}\pi
\]

当 $x$ 从 $0\to 1$ 时，$t$ 从 $\frac{\pi}{2}\to 0$，且 $\mathrm{d}x=-3\cos^{2}t\sin t\,\mathrm{d}t$，利用点火公式得：

\[
V_{\text{星}}=\pi\int_{\frac{\pi}{2}}^{0}\sin^{6}t(-3\cos^{2}t\sin t)\,\mathrm{d}t =3\pi\int_{0}^{\frac{\pi}{2}}\sin^{7}t(1-\sin^{2}t)\,\mathrm{d}t=3\pi\left(\frac{16}{35}-\frac{128}{315}\right)=\frac{16}{105}\pi
\]

\[
V = \frac{2}{3}\pi - \frac{16}{105}\pi = \frac{18}{35}\pi
\]

\textbf{3. 求表面积 $S = S_{\text{圆}} + S_{\text{星}}$}：

圆弧形成单位半球面，$S_{\text{圆}} = 2\pi$；星形线弧微分 $\mathrm{d}s=\sqrt{x'^{2}+y'^{2}}\,\mathrm{d}t=3\sin t\cos t\,\mathrm{d}t$（$0\le t\le\frac{\pi}{2}$），故：

\[
S_{\text{星}}=2\pi\int_{0}^{\frac{\pi}{2}}\sin^{3}t\cdot(3\sin t\cos t)\,\mathrm{d}t =6\pi\int_{0}^{\frac{\pi}{2}}\sin^{4}t\,\mathrm{d}(\sin t)=\frac{6}{5}\pi,\qquad S=2\pi+\frac{6}{5}\pi=\frac{16}{5}\pi
\]

\sechead{四、易错点}

\warnline{① 体积减、表面积加：最易把表面积也写成 $S_{\text{外}}-S_{\text{内}}$！挖去内部后，内壁恰恰增加了表面积，必须相加。}

\warnline{② 参数方程定限方向：星形线 $t=0$ 对应 $x=1$，$t=\frac{\pi}{2}$ 对应 $x=0$；用 $\mathrm{d}x=x'(t)\,\mathrm{d}t$ 时下限是 $\frac{\pi}{2}$、上限是 $0$，写反会得负体积。}

\warnline{③ 弧微分开方注意符号：$\mathrm{d}s=3\lvert\sin t\cos t\rvert\,\mathrm{d}t$，本题 $t\in[0,\frac{\pi}{2}]$ 才可直接去绝对值；跨象限必须分段讨论。}

\warnline{④ 勿混用体积与侧面积公式：体积系数为 $\pi$ 且积分 $y^{2}\,\mathrm{d}x$；侧面积系数为 $2\pi$ 且积分 $\lvert y\rvert\,\mathrm{d}s$（含根号弧微分）。}

\end{document}
